\documentclass{homework}
% \usepackage{graphicx} % Required for inserting images
\usepackage{amsmath}
\usepackage{gensymb}
\usepackage{subfig}
\usepackage{float}

\class{ECEN 5448}
\title{Homework 4}
\author{Sean Jennings}
\date{November 30, 2023}

\begin{document}
\maketitle

\section{Problem 1}
\subsection{}

\newcommand{\contr}{\mathcal{C}}

\begin{letterenum}
\item Take as a simple counter-example:
    \begin{equation*}
        A = \mat{
            0 & 1 \\
            1 & 0
        } \qquad B = \mat{
            0 \\ 1
        }.
    \end{equation*}
    Clearly, $\mathcal{C}(A, B) = \mat{0 & 1 \\ 1 & 0}$ is full rank, and
    thus $(A, B)$ is controllable. On the other hand, we have:
    \begin{equation*}
        A^2 = \mat{
            1 & 0 \\
            0 & 1
        }
    \end{equation*}
    so that the controllability matrix for the system $(A^2, B)$ is given by:
    \begin{equation*}
        \mathcal{C}(A^2, B) = \mat{B & A^2 B} = \mat{0 & 0 \\ 1 & 1}.
    \end{equation*}
    With a row of zeros, $\contr(A^2, B)$ is clearly not full-rank, so the system
    $(A^2, B)$ is not controllable, and thus the premise is disproven by
    counterexample.

\item The statement is true. Even with an additional $k(t)$ term in the system,
    the controllability of $(A, B)$ ensures that the inhomogenous system
    \begin{equation*}
        \dot{x} = Ax + Bu + k(t)
    \end{equation*}
    is also controllable. To see this, we modify the solution to the inhomogenous
    LTI system to include the $k(t)$ term:
    \begin{equation*}
        x(t) = \Phi(t, t_0) x(t_0) + \int_{t_0}^{t} \Phi(t, \tau) \bigl[
            B u(\tau) + k(t)
        \bigr] d\tau
    \end{equation*}
    where we have replaced $Bu(\tau)$ in the standard equation with
    $Bu(\tau) + k(\tau)$. $\Phi(t, t_0)$ is the state transition matrix.
    Extracting the $k(\tau)$ term from the integral yields:
    \begin{equation*}
        x(t) = \biggl[
            \Phi(t, t_0) x(t_0) + \int_{t_0}^t \Phi(t, \tau) k(\tau) d\tau
        \biggr] + \int_{t_0}^{t} \Phi(t, \tau) B u(\tau) d\tau.
    \end{equation*}

    Let $x_1 \in \R^n$ be an arbitrary state we'd like to reach at time $t_1 > t_0$.
    Subtract the term in the brackets from both sides of the equation, and 
    define the left-hand side as $\Delta x \in \R^n$:
    \begin{equation*}
        \Delta x := x_1 - \Phi(t_1, t_0) x(t_0) -
            \int_{t_0}^{t} \Phi(t, \tau) k(\tau) d\tau.
    \end{equation*}
    Since the system $(A, B)$ is controllable, the range of the grammian
    $R(W([t_0, t_1]))$ (where $R$ denotes the range) is all of $\R^n$, and we have
    $\Delta x \in R(W([t_0, t_1]))$.

    Let $v$ be the vector such that $\Delta x = W([t_0, t_1]) v$.
    Then, the controller $u(t)$
    \begin{equation*}
        u(t) = B^T \Phi^T(t, t_0) v
    \end{equation*}
    will bring the system from the arbitrary position $x(t_0)$ to any
    other position $x_1$ at time $t_1 > t_0$.

\item Yes, since the system is controllable, there exists a control $u$
    that brings the system to any $x_1 \in \R^n$ at any time $t_1$. Put
    $x_1 = 0_n$ and $t = 1$, and let $v$ be the vector that satisfies:
    \begin{equation*}
        v = W([0, 1])^{-1} (x(1) - \Phi(1, 0) x(0)).
    \end{equation*}
    (The grammian $W([0, 1])$ is invertible since it is a square matrix and
    controllability of $(A, B)$ implies it is full rank.)

    Then, the control
    \begin{equation*}
        u(t) = \biggl\{
            \begin{array}{ll}
                B^T \Phi(t, 0) v &\qquad \text{if } 0 \leq t < 1 \\
                0   &\qquad \text{otherwise}
            \end{array}
        \biggr.
    \end{equation*}
    will bring the system to rest at $t=1$, and it will remain at rest,
    since
    \begin{equation*}
        \dot x(t) = A 0_n + B 0_n = 0_n
    \end{equation*}
    for all $t \geq 1$.

\item While the system can be brought to any $\bar{x} \in \R^n$ at time $t=1$,
    it is not necessarily possible to maintain that state, even though the 
    system is controllable.

    To see this, put $x = \bar{x}$ and $\dot x = 0$ into the state equation:
    \begin{equation*}
        0 = A \bar{x} + Bu
    \end{equation*}
    This equation can only hold if $A \bar{x}$ lies in the range of $B$. 

    Take the following example:
    \begin{equation*}
        A = \mat{0 & 1 \\ -1 & 0} \qquad B = \bar x = \mat{0 \\ 1}.
    \end{equation*}
    Then the controllability matrix is given by:
    \begin{equation*}
        \mathcal{C}(A, B) = \mat{0 & 1 \\ 1 & 0}
    \end{equation*}
    and is clearly full rank. However,
    \begin{equation*}
        A \bar{x} = \mat{1 \\ 0} \notin R\biggl(\mat{0 \\ 1}\biggr) = R(\bar x)
    \end{equation*}
    so that $x = \bar x$ implies that $\dot{x} \neq 0$. So although
    $\bar{x}$ can be reached at even an infinite sequence of times
    $\{t_i \}_{i \in \N}$ with $t_{i+1} > t_i \geq 1$,
    it cannot be held there for all $t \geq 1$.

    The counter-example given could represent an undampened spring-mass system with 
    spring constant $k=1$. The result can then be explained in terms
    of mechanical intuition: while the system can be brought to any
    position at $t=1$, a target position $\bar x$ with non-zero velocity
    cannot be an equilibrium position. Regardless of the force applied, at the
    next instant, the mass will have moved.
\end{letterenum}

\subsection{}

\newcommand{\detr}{\text{det}}
We first compute the controllability matrix:
\begin{equation*}
    \mathcal{C}(A, B) = \mat{
        1 & 1 & 1 \\
        1 & -2 & 1 \\
        1 & -4 & 6
    }
\end{equation*}
Using a computer to find the determinant yields $\detr(\mathcal{C}(A, B)) = -15$.
Since the determinant is non-zero the matrix is invertible and the system
is therefore controllable.

To convert to controllable canonical form, given by
$(A', B') = (PAP^{-1}, PB)$, we begin by finding the characteristic polynomial:
\begin{align*}
    \detr(sI - A) &= \detr \biggl(
        \mat{
            s - 1 & 0 & 0 \\
            1 & s + 1 & 0 \\
            2 & 0 & s + 2 \\
        }
    \biggr) \\ 
    &= (s - 1)(s + 1)(s + 2) \\
    &= (s - 1)(s^2 + 3s + 2) \\
    &= s^3 + 2s^2 - s - 2.
\end{align*}
Then, $(A', B')$ can be immediately deduced from its coefficients
\begin{equation*}
    A' = \mat{
        0 & 1 & 0 \\
        0 & 0 & 1 \\
        2 & 1 & -2 \\
    } \qquad B' = \mat{0 \\ 0 \\ 1}
\end{equation*}
Since the control input is one dimensional, the matrix $P$
is given by:
\begin{align*}
    P &= \contr(A', B') \contr(A, B)^{-1} \\
        &= \mat{
            0 & 0 & 1 \\
            0 & 1 & -2 \\
            1 & -2 & 5 \\
        } \mat{
            1 & 1 & 1 \\
            1 & -2 & 1 \\
            1 & -4 & 6
        }^{-1} \\
        &= \frac{1}{15} \mat{
            2 & -5 & 3 \\
            1 & 5 & -6 \\
            8 & -5 & 12 \\
        }
\end{align*} 

\subsection{}

The controllability matrix is given by:
\begin{equation*}
    \contr(A, B) = \mat{1 & 0 \\ 1 & 0}.
\end{equation*}
So $v_1 = \mat{1 & 1}^T$ forms a basis for $R(\contr(A, B))$  and $v_2 = \mat{1 & -1}^T$ is one vector such that
$\{v_1, v_2\}$ spans $\R^2$.

Define $P$ such that $Pv_1 = e_1$ and $Pv_2 = e_2$ where $e_1,e_2$ are
the standard basis in $\R^2$. We can then find $P$ via the following 
system of equations:
\begin{align}
    P_{11} + P_{12} = 1 \label{eq:P1} \\
    P_{21} + P_{22} = 0 \label{eq:P2} \\
    P_{11} - P_{12} = 0 \label{eq:P3} \\
    P_{21} - P_{22} = 1 \label{eq:P4}.
\end{align}
where $P_{ij}$ is the $ij$th entry of $P$.
Adding (\ref{eq:P1}) and (\ref{eq:P3}), we obtain $P_{11} = P_{12} = 1/2$,
and similarly adding (\ref{eq:P2}) and (\ref{eq:P4}) yields
$P_{21} = 1/2$, $P_{22} = -1/2$. Then, $P$ is 
\begin{equation*}
    P = \frac{1}{2} \mat{1 & 1 \\ 1 & -1}.
\end{equation*}

Solving for $PAP^{-1}$ and $PB$ reveals the system in Kalman decomposition form:
\begin{align*}
    \dot{x}' &= PAP^{-1} x' + PBu \\
        &= \mat{0 & 0 \\ 0 & -2} x' + \mat{1 \\ 0}u
\end{align*}

\section{Observability of Linear Systems}
\subsection{}

Let $y(t)$ and $y'(t)$ be the output of the given LTV system with initial
conditions $x(t_0)$ and $x'(t_0)$, respectively. Then, substituting in
the solution for $x(t)$ and $x'(t)$ into the observation equation, we have:
\begin{align*}
    y(t) = C(t)x(t) &= C(t) \biggl[
        \Phi(t, t_0) x(t_0) + \int_{t_0}^{t} \Phi(t, \tau) B(\tau) u(\tau) d\tau
    \biggr] \\
    y'(t) &= C(t) \biggl[
        \Phi(t, t_0) x'(t_0) + \int_{t_0}^{t} \Phi(t, \tau) B(\tau) u(\tau) d\tau
    \biggr]
\end{align*}
Then, subtracting these two equations, we find
\begin{equation}
    y(t) - y'(t) = C(t) \Phi(t, t_0) \biggl[
        x(t_0) - x'(t_0)
    \biggr]
    \label{eq:output-eq}
\end{equation}
and therefore $x(t_0)$ and $x'(t_0)$ are output equivalent if and only if
equation (\ref{eq:output-eq}) is equal to 0. 

($\Rightarrow$) Now, assume that the initial conditions are output equivalent and 
multiply the observability grammian $M(t_0, t_1) \in \R^{n\times n}$
by the difference in initial conditions:
\begin{align*}
    M(t_0, t_1) (x(t_0) - x'(t_0)) &= \biggl(
        \int_{t_0}^{t_1} \Phi^T(\tau, t_0) C^T(\tau) C(\tau) \Phi(\tau, t_0) d\tau
    \biggr) (x(t_0) - x'(t_0)) \\
    &= \int_{t_0}^{t_1} \Phi^T(\tau, t_0) C^T(\tau) \biggl[
        C(\tau) \Phi(\tau, t_0) (x(t_0) - x'(t_0))
    \biggr]d\tau \\
    &= \int_{t_0}^{t_1} \Phi^T(\tau, t_0) C^T(\tau) (y(\tau) - y'(\tau)) d\tau \\
    &= 0.
\end{align*}
Therefore the difference $(x(t_0) - x'(t_0))$ is in the null space of the grammian.

\newcommand{\nul}{\mathcal{N}}

($\Leftarrow$) Suppose $x(t_0) - x(t_1) \in \nul(M(t_0, t_1))$.
For any $\eta\in\nul(M(t_0, t_1))$, we have $M(t_0, t_1)\eta = 0_n$ and $\eta^T M(t_0, t_1) \eta = 0$.
So then,
\begin{align*}
    \eta^T M(t_0, t_1) \eta &= \eta \biggl(
        \int_{t_0}^{t_1} \Phi^T(\tau, t_0) C^T(\tau)
            C(\tau) \Phi(t, \tau) d\tau
    \biggr) \eta \\
    &= \int_{t_0}^{t_1} v^T(\tau) v(\tau) d\tau \\
    &= \int_{t_0}^{t_1} \norm{v(\tau)}^2 d\tau \\
    &= 0
\end{align*}
where we have defined $v(\tau) := C(\tau) \Phi(\tau, t_0) \eta$.
The norm is a positive definite function.
Therefore, for the integral to be zero, we must have $v(\tau) = 0$ for all
$\tau \in [t_0, t_1]$. Putting $\eta = x(t_0) - x'(t_0)$  and $\tau=t$,
we have:
\begin{equation*}
    C(t) \Phi(t, t_0) \biggl[
        x(t_0) - x'(t_0)
    \biggr] = 0
\end{equation*}
and by equation (\ref{eq:output-eq}), the initial conditions are 
output equivalent.

\subsection{}

Let $A_1, A_2 \in\R^{3\times 3}$ be the piecewise constant values of
$A(t)$ during the first and second time intervals, respectively. We begin by finding
the jordan form $A_1 = P_1 D_1 P_1^{-1}$ for the first matrix. A computer
outputs:
\begin{equation*}
    P_1 = \frac{1}{\sqrt{2}} \mat{
        1 & -j & 0 \\
        1 & j & 0 \\
        0 & 0 & \sqrt{2} \\
    } \qquad D_1 = \mat{
        j & 0 & 0 \\
        0 & -j & 0 \\
        0 & 0 & 0 \\
    }
\end{equation*}
where $j := \sqrt{-1}$.

Then, we can find the state transition matrix $\Phi_1(t, 0)$, which is valid
on the interval $[0, \pi)$:
\begin{align*}
    \Phi_1(t, 0) &= P_1 e^{D_1t} P_1^{-1} \\
        &= \frac{1}{\sqrt{2}} \mat{
            1 & 1 & 0 \\
            j & -j & 0 \\
            0 & 0 & \sqrt{2}
        } \mat{
            e^{jt} & 0 & 0 \\
            0 & e^{-jt} & 0 \\
            0 & 0 & 1 \\
        } P_1^{-1} \\
        &= \frac{1}{\sqrt{2}} \mat{
            e^{jt} & e^{-jt} & 0 \\
            je^{jt} & -je^{-jt} & 0 \\
            0 & 0 & \sqrt{2} \\
        } \frac{1}{\sqrt{2}} \mat{
            1 & -j & 0 \\
            1 & j & 0 \\
            0 & 0 & \sqrt{2} \\
        } \\
        &= \frac{1}{2} \mat{
            e^{jt} + e^{-jt} & j (e^{-jt} - e^{jt}) & 0 \\
            -j (e^{-jt} - e^{jt}) & e^{jt} + e^{-jt} & 0 \\
            0 & 0 & 2
        } \\
        &= \mat{
            \cos(t) & \sin(t) & 0 \\
            -\sin(t) & \cos(t) & 0 \\
            0 & 0 & 1
        }
\end{align*}

where we have utilized the Euler-related identities:
\begin{align*}
    \frac{1}{2} (e^{jt} + e^{-jt}) = \cos(t) \\
    \frac{j}{2} (e^{-jt} + e^{-jt}) = \sin(t) \\
\end{align*}

We have that
\begin{equation*}
    c \Phi_1(t, 0) = \mat{-\sin(t) & \cos(t) & 0}
\end{equation*}

and thus the grammian over any subset of the first interval
$[t_0, t_1] \subseteq [0, \pi]$ is given by:
\begin{align*}
    M(t_0, t_1) &= \int_{0}^{t} (\Phi_1^T(\tau, 0) c^T) (c \Phi_1(\tau, 0)) d\tau \\
        &= \int_{0}^{t} \mat{
            -\sin(\tau) \\
            \cos(\tau) \\
            0
        } \mat{-\sin(\tau) & \cos(\tau) & 0} d\tau \\
        &= \int_{0}^{t} \mat{
            \sin^2(\tau) & -\cos(\tau) \sin(\tau) & 0 \\
            -\cos(\tau) \sin(\tau) & \cos^2(\tau) & 0 \\
            0 & 0 & 0
        } d\tau.
\end{align*}

Clearly, the observability grammian is not full rank, and
thus the system is not observable over this interval. 

The computation for the state transition matrix $\Phi_2(t, \pi)$ over the second
interval $t \in [\pi, 2\pi]$ follows very similarly to the first and can be
worked out to be:
\begin{align*}
    \Phi_2(t, \pi) &= e^{A_2(t - \pi)} \\
        &= \mat{
            1 & 0 & 0 \\
            0 & \cos(t - \pi) & \sin(t - \pi) \\
            0 & -\sin(t - \pi) & \cos(t - \pi) \\
        }.
\end{align*}
So,
\begin{equation*}
    c \Phi_2(t, \pi) = \mat{0 & \cos(t - \pi) & \sin(t - \pi)}
\end{equation*}
and the grammian over any subset $[t_0, t_1] \subseteq [\pi, 2\pi]$
is given by
\begin{equation*}
    M(t_0, t_1) = \int_{t_0}^{t_1} \mat{
        0 & 0 & 0 \\
        0 & \cos^2(t - \pi) & \cos(t - \pi) \sin(t - \pi) \\
        0 & \cos(t - \pi) \sin(t - \pi) & \sin^2(t - \pi)
    }.
\end{equation*}
Again, the observability grammian is clearly singular and the system
is not observable over $[\pi, 2\pi]$.

However, over the interval $[0, 2\pi]$ the grammian is given by:
\begin{align*}
    M(0, 2\pi) &= M(0, \pi) + M(\pi, 2\pi)  \\
        &= \int_{0}^{\pi} \mat{
            \sin^2(\tau) & -\cos(\tau) \sin(\tau) & 0 \\
            -\cos(\tau) \sin(\tau) & \cos^2(\tau) & 0 \\
            0 & 0 & 0
        } \\
        &\quad + \int_{\pi}^{2\pi} \mat{
            0 & 0 & 0 \\
            0 & \cos^2(t - \pi) & \cos(t - \pi) \sin(t - \pi) \\
            0 & \cos(t - \pi) \sin(t - \pi) & \sin^2(t - \pi)
        } \\
        &= \frac{\pi}{2} \mat{
            1 & 0 & 0 \\
            0 & 2 & 0 \\ 
            0 & 0 & 1
        }
\end{align*}
where the following equalities (which follow from double angle identities)
were used
\begin{align*}
    \int_{0}^{\pi} \cos^2(t) dt = 
    \int_{\pi}^{2\pi} \cos^2(t - \pi) dt &=  \frac{\pi}{2} \\
    \int_{0}^{\pi} \sin^2(t) dt = 
    \int_{\pi}^{2\pi} \sin^2(t - \pi) dt &= \frac{\pi}{2} \\
    \int_{0}^{\pi} \cos(t) \sin(t) dt = 
    \int_{\pi}^{2\pi} \cos(t - \pi) \sin(t - \pi) dt &= 0.
\end{align*}

Thus, the grammian is non-singular and the system is observable over the full
interval $[0, 2\pi]$ (and quite possibly any subset of $[0, 2\pi]$ which intersects
both $[0, \pi]$ and $[\pi, 2\pi]$, although this makes the math a fair bit messier).

\section{Transfer Function and Realization}
\subsection{}
We rewrite the system in matrix form:
\begin{align*}
    \dot{x} &= \mat{
        0 & 1 \\
        -3 & -2
    } x + \mat{-1 \\ 1} u \\
    y &= \mat{1 & 1} x.
\end{align*}

Since the system is SISO, taking the Laplace transfer yields:
\begin{equation*}
    \hat(y) C(sI - A)^{-1} B \hat{u}
\end{equation*}
where $\hat{y}(s) = \mathcal{L}(y(t))$ and $\hat{u}(s) = \mathcal{L}(u(t))$
are the Laplace domain representations of time-domain signals $y$ and $u$.

We focus first on computing $(sI - A)^{-1}$.
\begin{align*}
    sI - A = \mat{
        s & -1 \\
        3 & s + 2 \\
    } \\
    \detr(sI - a) = s^2 + 2s + 3 \\
    (sI - a)^{-1} = \frac{1}{s^2 + 2s + 3} \mat{
        s + 2 & 1 \\
        -3 & s 
    } \\
\end{align*}
Then, the transfer function $G(s)$ for the system can be computed:
\begin{align*}
    G(s) = C(sI - A)^{-1} B &= \frac{1}{s^2 + 2s + 3} \mat{1 & 1} \mat{
        s + 2 & 1 \\
        -3 & s 
    } \mat{-1 \\ 1} \\
    &= \frac{1}{s^2 + 2s + 3} \mat{s - 1 & s + 1} \mat{-1 \\ 1} \\
    &= \frac{2}{s^2 + 2s + 3}.
\end{align*}

\subsection{}

To find a minimum realization, we find coefficients of the numerator
and denominator:
\begin{align*}
    p(s) = p_2 s^2 + p_1 s + p_0 \\
    p_0 = 1 \qquad p_1 = 3 \qquad p_2 = 2 \\
    q(s) = s^3 + q_2 s^2 + q_1 s + s_0 \\
    q_0 = 1 \qquad q_1 = -3 \qquad q_2 = -3
\end{align*}

The negative of the coefficients of $q$ become the final row of
the CCF of $A$, and the coefficients of $p$ become the observability
row vector $C$:
\begin{align*}
    A &= \mat{
        0 & 1 & 0 \\
        0 & 0 & 1 \\
        -1 & 3 & 3 \\
    } &\qquad B &= \mat{0 \\ 0 \\ 1} \\
    C &= \mat{1 & 3 & 2} &\qquad D &= 0
\end{align*}

At last, $(A', B', C', D') = (A^T, C^T, B^T, D)$ is the dual system,
and by duality, will be in OCF:
\begin{align*}
    A' &= \mat{
        0 & 0 & -1 \\
        1 & 0 & 3 \\
        0 & 1 & 3 \\
    } &\qquad B' &= \mat{1 \\ 3 \\ 2} \\
    C' &= \mat{0 & 0 & 1} &\qquad D' &= 0
\end{align*}



















\end{document}
